\documentclass[a4paper]{report}
% \usepackage{predefine}

% 中文支持和 Unicode 编译
\usepackage[UTF8]{ctex}
% 数学支持
\usepackage{amsmath,amssymb,mathtools,amsthm}
% 绘图支持
\usepackage{tikz}
\usetikzlibrary{arrows.meta,decorations.pathmorphing,decorations.pathreplacing,patterns,angles,quotes}

% 页面和超链接
\usepackage{geometry}
\usepackage[colorlinks=true,linkcolor=blue,citecolor=blue,urlcolor=blue]{hyperref}
\usepackage{cite}

% 尺寸
\geometry{a4paper,top=30mm,bottom=30mm,left=30mm,right=30mm}
\setlength{\parskip}{0.8em}
\setlength{\parindent}{2em}

% 缩写
\newcommand{\action}{\mathcal{S}}
\newcommand{\lag}{\mathcal{L}}
\newcommand{\ham}{\mathcal{H}}
\newcommand{\bd}{\mathrm{d}}

\newtheorem{theorem}{定理}[section]



\author{Cheng Li}
\title{\textbf{{\Huge 经典场论}}\\{\normalsize 连续介质提力学与电动力学}}
\begin{document}
\maketitle
\tableofcontents

\part{连续介质体力学}
\chapter{弹性体力学}
\section{位移场}
\section{应力与应变张量}
Navi-Cauchy
\begin{equation}
    (\lambda+\mu) \nabla(\nabla\cdot \uf) + \mu \nabla^2 \uf + \boldsymbol{f} = \rho \frac{\partial^2\uf}{\partial t^2}
\end{equation}
\section{弹性波}
\chapter{流体力学}
\section{流速场}
连续性方程
\begin{equation}
    \frac{\partial\rho}{\partial t}+ \nabla\cdot(\rho\uf)=0
\end{equation}
\section{流体静力学}
\begin{equation}
    \nabla p = \rho \boldsymbol{f}
\end{equation}
\section{理想流体动力学}
\begin{equation}
    \nabla\cdot\uf = 0
\end{equation}
Euler-equations
\begin{equation}
    \rho \left( \frac{\partial\uf}{\partial t}+ (\uf\cdot\nabla)\uf \right) = -\nabla p + \rho\boldsymbol{f}
\end{equation}
\section{粘性流体动力学}
Navi-Stokes equations
\begin{equation}
    \rho \left( \frac{\partial\uf}{\partial t}+ (\uf\cdot\nabla)\uf \right) = -\nabla p + \mu\nabla^2 \uf+ \rho\boldsymbol{f}
\end{equation}
\section{空气动力学}

\part{电动力学}
Gaussian Units
\begin{equation}
    \varepsilon_0 = \frac{1}{4\pi}, \qquad \mu_0 = \frac{4\pi}{c^2}
\end{equation}
\chapter{电磁场基本理论}
\section{麦克斯韦方程组}
\begin{equation}
    \begin{split}
        \nabla\cdot \boldsymbol{E} &= 4\pi{\rho_c}\\
        \nabla\times\boldsymbol{E} &= -\frac{1}{c}\frac{\partial \boldsymbol{B}}{\partial t}\\
        \nabla\cdot\boldsymbol{B}&=0\\
        \nabla\times\boldsymbol{B}&=\frac{4\pi}{c}\boldsymbol{J}+\frac{1}{c}\frac{\partial\boldsymbol{E}}{\partial t}
    \end{split}
\end{equation}

\begin{equation}
    \action = \int \bd t~\bd^3\boldsymbol{x}~\frac{1}{2}\Big[ (\nabla\phi)^2-(\nabla\times\boldsymbol{A})^2  \Big]
\end{equation}
\section{极化与磁化}
\section{电磁场边界条件}
\chapter{静电场}
\section{电标势}
\section{电偶极子}
\section{连续电荷系统及电偶极展开}
\section{导体的静电场}
\chapter{静磁场}
\section{库伦规范与磁矢势}
\section{磁偶极子}
\section{铁磁体}
\chapter{电磁波}
\section{洛伦兹规范与达朗贝尔方程}
\section{电磁辐射}
% \section{电磁波的传播}
\chapter{动体的电动力学}
\section{运动带电粒子与电磁场的相互作用}
\section{运动带电粒子的辐射}
% \chapter{协变形式的电磁场理论}
% \section{闵可夫斯基空间与四维协变量}
% \section{电磁场张量}
\appendix
\chapter{矢量分析}
\chapter{偏微分方程与特殊函数}


\end{document}